\subsection{Propagation between components in degrees two and three}

In the ordinary deformation argument the closed set to be propagated
is \(V(J)\), the inverse image of the specified Hecke ideal.
Here are the component and curve arguments with their dimension and
point hypotheses visible.

\begin{lemma}\label{cprop:curve}
Let \((R,\mathfrak n)\) be a Noetherian local ring, \(\pi\in\mathfrak n\),
\(b\ge1\), and \(I,H_1,\ldots,H_t\) ideals. Suppose
\(\dim R/I\ge b+1\) and \(\dim R/H_i<b\) for each \(i\).
There is a prime \(\mathfrak p\supset I+(\pi)\) with
\(\dim R/\mathfrak p=1\) containing none of the \(H_i\).
\end{lemma}
\begin{proof}
In \(R/I\) a system of parameters modulo \(\pi\), together with
\(\pi\), generates an ideal primary to the maximal ideal.
Krull's height theorem gives \(\dim R/(I+(\pi))\ge b\). Take a
minimal prime \(Q\) of this quotient of dimension at least \(b\),
viewed as a prime of \(R\). Each \(H_i\) contains an element
\(a_i\notin Q\), by the strict inequality of dimensions.
Choose \(u\in\mathfrak n\setminus Q\) and put \(z=u\prod_i a_i\).
Choose a prime \(\mathfrak p\) maximal among primes containing \(Q\)
and avoiding \(z\). In the local domain \(D=R/\mathfrak p\), the
image of \(z\) is a nonzero nonunit belonging to every nonzero prime.
If \(\dim D\ge2\), the height-one primes containing \(z\) are the
finitely many minimal primes over \((z)\), by the principal ideal
theorem. Prime avoidance gives a nonunit in the maximal ideal
outside all of them. A minimal prime over that nonzero element
has height one and also contains \(z\), a contradiction. Thus
\(\dim D=1\), and the choice of \(z\) excludes every \(H_i\).
\end{proof}
For the source chart of rank \(n=2,3\) and layer degree parameter
\(d\), take \(b=nd\). The forbidden ideals are the opposite triangular
ideal and, for \(i\ne j\) at each marked \(v\), the sum of the Hecke
ideal, the coefficient maximal ideal and the ideal of all differences
\(\psi_{vi}(g)-\psi_{vj}(g)\) on the specified inertia subgroup.
If the first triangular ideal already generates the unit ideal on
the chart and each forbidden quotient has dimension \(<nd\), the
lemma produces a characteristic-\(\ell\) curve off both triangular
loci, with all pairs of marked characters distinct there. For
the difference ideals this uses that the initial prime contains
the Hecke ideal and that the curve contains \(\pi\).
\src{src/Development/AutomorphicArithmetic/AutomorphicVarieties/CurveDVRSpecialization/NicePrimeEx.lean}{87}{The one-dimensional prime} and
\src{src/Development/AutomorphicArithmetic/AutomorphicVarieties/CurveDVRSpecialization/ExistsNicePrimeAboveOfLargeDim.lean}{7}{the exceptional ideals in the chart}.

\paragraph{The path through a local chart.}
Write \(U\) for the open locus off both triangular ideals and, in
rank three, off the ideal of matrix entries of odd weight difference.
Write \(E_0\) for the part off both triangular ideals where that
third ideal does vanish. In degree two \(E_0\) is empty. Call a prime
\emph{selected} when it is a one-dimensional-quotient prime as just
described (including the separation and characteristic conditions).
It has type \(U\) or \(E_0\). In the next lemma a designated subset
of component primes of a chart ideal is allowed to carry the
exceptional links of the rank-three construction. Set \(D=nd+1\)
and \(d(P)=\dim R/P\).

\begin{lemma}\label{cprop:local}
Let \(J\) define the closed Hecke locus in a Noetherian chart.
Let \(\mathcal P\) be the minimal primes of a specified chart ideal,
including every minimal prime of \(R\) lying in \(U\). Assume each
vertex of \(\mathcal P\) is either in \(U\) or designated exceptional.
Join \(P,Q\in\mathcal P\) when \(\dim R/(P+Q)\ge D\), and assume
this graph is connected. Suppose the following point assertions hold.
\begin{enumerate}
\item Every prime over \(J\) of quotient dimension at least \(D\)
has a selected specialization. Above any designated exceptional
vertex, every prime of that dimension has a selected specialization
of type \(E_0\).
\item At a selected \(p\in U\) with \(J\subset p\), every prime
below \(p\) contains \(J\).
\item Link selected \(E_0\)-primes when they share a designated
exceptional vertex below them. If \(p_0\) has a \(U\)-prime
containing \(J\) below it, every \(U\)-prime below every selected
\(E_0\)-prime linked by a finite path to \(p_0\) contains \(J\).
Paths of length zero are included.
\end{enumerate}
If one \(U\)-prime \(I\) contains \(J\) and \(d(I)\ge D\), every
minimal prime of \(R\) in \(U\) contains \(J\).
\end{lemma}
\begin{proof}
At any selected \(p\) above a \(U\)-prime containing \(J\), all
\(U\)-primes below \(p\) contain \(J\): use assertion 2 or the
length-zero case of assertion 3. Choose a minimal prime \(P_0\)
below \(I\); openness gives \(P_0\in U\). Applying assertion 1
at \(I\) and the preceding observation gives \(J\subset P_0\).

Travel along a path in \(\mathcal P\) from \(P_0\). At a \(U\)-vertex
retain the conclusion that it contains \(J\). At any vertex outside
\(U\) retain a selected \(E_0\)-prime above it reachable from a
selected \(p_0\in E_0\) above some \(U\)-prime containing \(J\).
For an edge \(P,Q\) take a prime \(M\supset P+Q\) with \(d(M)\ge D\),
from a maximal-dimensional component of \(R/(P+Q)\).
If \(Q\in U\) and \(P\in U\), assertion 1 at \(M\supset J\)
and the observation give \(J\subset Q\).
If \(Q\in U\) and \(P\notin U\), select an \(E_0\)-prime above
\(M\) by assertion 1 at the designated vertex \(P\). The retained
path extends to it through \(P\). Assertion 3 gives \(J\subset Q\).
If \(Q\notin U\), select an \(E_0\)-prime above \(M\) using \(Q\).
For \(P\in U\) it starts a path with witness \(P\). For \(P\notin U\)
it extends the retained path through \(P\). This proves the claimed
induction. Every minimal prime in \(U\) is a vertex of the connected graph.
\end{proof}
\src{src/Development/AutomorphicArithmetic/AutomorphicVarieties/CurveDVRSpecialization/LocProp.lean}{17}{Local propagation}
and \src{src/Development/AutomorphicArithmetic/AutomorphicVarieties/CurveDVRSpecialization/PropagationTheoremGL2.lean}{32}{the two prime-transfer hypotheses}.

\begin{lemma}\label{cprop:global}
Let \(X\) be Noetherian, \(U\subset X\) open and \(V(J)\) closed.
Call an irreducible component eligible if it meets \(U\).
Let \(B\subset X\) be connected and let \(T\subset B\) be test points
such that each closed set meeting \(B\) meets \(T\).
Suppose every eligible component meets \(B\), their union contains
a dense subset of \(B\), and at any \(t\in T\), containment in
\(V(J)\) passes between eligible components through \(t\).
If one eligible component lies in \(V(J)\), then \(U\subset V(J)\).
\end{lemma}
\begin{proof}
Form the unions of the eligible components contained in \(V(J)\)
and of the others. Both unions are closed by finiteness. Together
they cover \(B\) by the density assumption. If the second group
exists, both groups meet \(B\), so connectedness gives an intersection
over \(B\). The intersection is closed, hence contains a point of
\(T\). The asserted transfer there contradicts membership of a
component in the second group. Every point of \(U\) lies on an
eligible component, proving the result.
\end{proof}

For clarity, transfer at \(t\) follows from Lemma~\ref{cprop:local}
on a flat local chart if each eligible component through \(t\)
has a chart minimal prime mapping to its generic point, of quotient
dimension at least \(D\). Indeed openness pulls the generic point
into \(U\). The first component supplies a \(J\)-containing prime,
the local lemma puts the other preimage prime over \(J\), and closure
gives the other whole component. For an affine flat local chart
centered at \(t\) the required lifting of generic points follows by
going down from its closed point to each minimal prime under \(t\),
and then taking a minimal prime below the lift. The contraction is
still the same minimal prime. The dimension inequality is additional.
In the global argument \(B\) is the connected locus used to test the
deformations, \(T\) its specified maximal points, and the density
required here is obtained in the source away from the closed
odd-weight-entry locus. When \(n=3\) the exceptional vertices also
use a finite ideal obtained from degree two. If the Hecke ideal \(J\)
works on \(U\) and the exceptional ideal \(I_3\) on \(E_0\),
then \(J\cap I_3\) works on their union. Finiteness of its fibre follows
from the radical test: any prime over a sum of the weight maximal ideal
and \(J\cap I_3\) contains \(J\) or \(I_3\), since \(JI_3\subset J\cap I_3\).
Thus the closed fibre in the Noetherian local ring is zero-dimensional,
hence Artinian. Its composition factors are its residue field, finite
over the coefficient residue field in the ordinary deformation chart.
This gives finite dimension over that coefficient field.
This makes explicit the use of both ranks. It still requires the
dimension and connectedness inputs, a seed on the actual tower and
the two arithmetic point-transfer assertions in the local lemma.
\src{src/Development/AutomorphicArithmetic/AutomorphicVarieties/CurveDVRSpecialization/GlobProp.lean}{3}{Global propagation},
\src{src/Development/AutomorphicArithmetic/AutomorphicVarieties/CurveDVRSpecialization/GlobOfCharts.lean}{4}{passage through flat charts} and
\src{src/Development/AutomorphicArithmetic/AutomorphicVarieties/CurveDVRSpecialization/AbsMain.lean}{12}{the order of the two rank arguments}.
